% ============================================================
% content/answers-short.tex —— 2026 湖南卷物理 · 简明答案速查
% ============================================================

\begin{gknotice}
  选择题答案：1--B，2--D，3--C，4--B，5--A，6--A，7--D，8--BC，9--BC，10--AD
\end{gknotice}

\section*{一、选择题}

\noindent
\begin{tabular}{@{}llllllllll@{}}
\toprule
题号 & 1 & 2 & 3 & 4 & 5 & 6 & 7 & 8 & 9\\
\midrule
答案 & B & D & C & B & A & A & D & BC & BC\\
\bottomrule
\end{tabular}

\smallskip

\noindent
第 10 题：AD

\section*{二、实验题}

\noindent
\textbf{11.} (3) $0.6\pi$~(s)；\quad(4) $b$；\quad(5) $g=\dfrac{4\pi^{2}l}{T^{2}}$；\quad(6) 高于

\smallskip

\noindent
\textbf{12.} (2) $c$、$d$；\quad50~mV；\quad(3) 11.350~mm（11.349--11.351 均可）；\quad(5) 见解析图像；\quad(6) 1.55~mm

\section*{三、解答题}

\noindent
\textbf{13.} (1) $n=\dfrac{2\sqrt{3}}{3}$；\quad(2) $\sin\theta=0.9$

\smallskip

\noindent
\textbf{14.} (1) $a_{0}=\dfrac{BEL}{mr}$；\quad(2) $v=\dfrac{E}{BL}-\dfrac{BLU_{0}}{mk}$；\quad(3) $Q=\dfrac{(R+2r)mE^{2}}{2(R+r)B^{2}L^{2}}$

\smallskip

\noindent
\textbf{15.} (1) $v=\sqrt{2gL}$；\quad(2) $\tan\alpha=\dfrac{\sqrt{23}}{2}$；\quad(3)(i) $E=\left[\dfrac{(1+k)^{2}}{4k\cos^{2}\theta}-1\right]mgL$；
\qquad(3)(ii) $d=\dfrac{8}{9}\sqrt{69}\,L\tan\theta$
